% Page budget: 1.55 pages | Owns: augmentation topology, additional states, encoder, and closed-loop objective.
% WRITING GUIDE
% Purpose: Define the final augmentation architecture and explain how it is trained inside the known feedback loop.
% Start here: Current implementation decisions D-129, D-140, D-141, D-180, D-220, D-233, and D-237 in docs/decisions.md.
% Supporting material: Quinten's 18 September outline feedback, the Meeting-audit synthesis, Thesis-writeup/Documentation/TRAINING-DESIGN.md, and docs/writeup figures.
% Mandatory papers: Read Hoekstra's 2025 LFR augmentation paper, the 2026 first-principles augmentation paper, the encoder-initialization paper, and Kessels' Chapter 5 before writing this section.
% Research-plan anchor: Preserve Aspect 2's dynamic parallel augmentation objective; document the departure from open-loop training and earlier state routing.
% Current decisions: Separate physical and additional states, keep the controller outside the model, score the complete training window without burn-in, and use one network for the physical correction and added-state update.
% Choices to justify: Final reported routing, added-state dimension, ANN width and depth, zero initialization, encoder history, closed-loop objective, rollout length, full-window scoring, and validation horizon. Use the configuration of the reported thesis run, not a superseded routing experiment.
% Horizon check: Treat encoder history, scored training window, joint-estimation horizon, and free-run validation horizon separately; relate candidates to relevant stable modes and test sensitivity to slower dynamics and free-integrator accumulation.
% Implementation audit: Read scripts/gantry/gantry_interconnect_dynamic.py, gantry_dynamic/{model,config,controller,data,training}.py, and model_augmentation/fit_systems/{interconnect,closed_loop,pre_encoder}.py.
% Reference check: Verify sources for parallel LFR augmentation, encoder initialization, closed-loop state extension, and any claimed architectural precedent against the original paper or thesis.
% Cautions: Earlier routing plans are superseded; distinguish the truth-only hidden absorber from learned states; closed-loop fit alone does not prove autonomous plant recovery.
% Figures: Absorber schematic, author's updated architecture, and thesis-owned common-reference loops. No residual panel or horizon figure.
% Section is finished when: Every signal has a defined frame and timing, the RK4 boundary is explicit, and the described architecture matches the final code.
%
% SOURCE MAP
% Hoekstra et al. 2025 EJC: dynamic-parallel topology, additional states, encoder-based truncated simulation.
% Hoekstra et al. 2026 Automatica submission: extended LFR formulation, baseline-preserving model/encoder initialization, joint identification context.
% Hoekstra et al. 2026 encoder paper: reconstructability-based W^b initialization; it does not identify W^a or test dynamic augmentation.
% Kessels 2025 thesis: direct closed-loop rollout, controller equations, and reconstruction of the machine controller state for each window.
% Drenth 2025 thesis: LPV-specific augmentation precedent; not the unrestricted MLP structure used here.
% This project: gantry specialization, residual-loop implementation, exact routing, RK4 boundary, and verification.
%
% MUST ESTABLISH (agreed with Dirk, 2026-10-06; step 2a of the thesis-section skill)
% Contribution delivered: the gantry realisation of S-DP augmentation (CT self-scheduled LPV-LFR baseline
%   inside a DT parallel model, joint estimation in the ten combinations) and its identification under
%   feedback. Not new: the S-DP topology (Hoekstra), augmentation of self-scheduled LPV-LFR with added
%   states (Drenth Eq. 5.2), closed-loop training with a known controller (Kessels, inspiration only).
% III-A Augmented model
%  1. Dynamic, not static: static augmentation adds no states (EJC Sec. 2, "flexible modes"); refitting
%     theta_base adds none either. Argue from the general case, no absorber.
%  1b. The augmented model is Hoekstra's LFR augmentation structure (EJC Eq. 4, 2026 Eq. 6) with a FIXED
%     interconnection (EJC Sec. 3.2) realising S-DP; its baseline block is Section II's self-scheduled
%     LFR with Delta(Y) = Y I6: an LFR inside an LFR, which is what the section title means.
%     Well-posedness: acyclic graph (EJC Thm. 1, 2026 Cond. 6) plus Section II's M(Y) condition.
%  2. Parallel, not series. LEAD REASON (Dirk, 2026-10-06): the orthogonal-by-construction method that
%     Section IV builds on is defined for the additive (parallel) structure (Gyorok 2026 abstract, Sec. 1;
%     D-003), so parallel lets us use that existing work. Supporting: Hoekstra gives no universal rule,
%     the structure follows prior knowledge (2026 Sec. 1, 6.3); the baseline stays the explicit nominal
%     term (research plan Aspect 2); parallel needs only a feedforward network for baseline-preserving
%     init, series needs ResNets (EJC Sec. 4.3, 2026 Sec. 5.4, 6.3).
%  3. State augmentation only, h_aug = 0: output P^T q is exact kinematics (cf. Kessels Remark 5.3). TODO reason.
%  4. Correction on all physical rows and added states: S-DP definition (EJC), X and Y reachable (D-103);
%     marginal-mode risk on free axes stated. TODO: why correct position rows rather than keep kinematics.
%     Context: Hoekstra 2026 Sec. 7 (F1Tenth) detaches the free integrators and learns only the
%     input-to-velocity map because the full state is measured; here only positions are measured.
%     One network for f_aug and g_aug is an implementation statement (D-180), no scientific claim.
%  5. DT correction after the CT RK4 step (input held, M(Y) per stage, affine normalisation). TODO reason
%     versus force injection in the CT derivative.
%  6. Unrestricted nonlinear augmentation (depends on Y through the state) versus LPV-structured
%     augmentation (Drenth): departure from the research plan's "in LPV-LFR form". TODO reason. Acyclic
%     routing adds no instantaneous loops; baseline well-posedness is Section II's.
%  7. Zero output layer reproduces the baseline from the same physical initial state (Hoekstra 2026
%     Eq. 29, D-237); the split is not unique (Sec. 5.2), handed to Section IV. Added states are memory
%     and model order, defined up to a state transformation, no physical coordinates.
% III-B Encoder
%  8. Encoder needed: only positions measured; truncated windows for cost and stability (Beintema).
%  9. W^b from the reconstructability map of a ZOH linearisation at an equilibrium (encoder paper
%     Eqs. 15-17, 30-32), not fitted (Hoekstra 2026 Eq. 30); W^a random. Claim only better starting
%     values and faster convergence, final accuracy comparable (encoder paper Sec. 4.4). The paper
%     assumes a stable baseline; ours is marginally stable (free X, Y), the map exists because the
%     positions make the linearisation observable. The paper treats co-estimating theta_base as out of
%     scope (footnote 1); we co-estimate. Precedent for linearising at zero: encoder paper Sec. 4.2
%     ("around the 0 point"), no reason given. TODOs: why Y = 0; disclose nominal parameters in W^b
%     versus detuned start; why model-based rather than data-based init (candidate: cost, Sec. 4.4);
%     Kaiming versus Xavier (todo only).
% 10. Encoder estimates the full state incl. measured positions (SUBNET). TODO: mention Kessels Remark 5.3?
% 11. Keep encoder history, scored window and validation horizon apart; values in Setup. The
%     joint-estimation sensitivity horizon belongs to Section IV.
% III-C Closed-loop identification
% 12. Known-controller, reference-driven output-error identification, inspired by Kessels (Eqs. 5.12,
%     5.13; Remark 5.6 open-loop attempt failed). Related to the indirect approach (Forssell and Ljung
%     Sec. 5.2) as a linear-theory motivation, not a guarantee; direct identification is not called
%     inherently biased. Free X and Y axes: the SUBNET consistency argument that Hoekstra inherits
%     (2026 Sec. 5.5) assumes an incrementally exponentially output stable data-generating system
%     (Beintema 2023 Condition 1); the open-loop gantry is not (free integrators), the loop with the
%     stabilising controller can be. Precondition, not a proof, for our nonlinear model. Same principle
%     as Kessels (closed-loop simulation of an unstable open-loop system), scoped to marginal stability.
% 13. Our formulation (ours, "we"): u_hat = u + C x^c + D(y - y_hat). Derivation by subtracting the
%     machine's controller equations, equal to the shared-r, shared-u_ff loop of the figure under a
%     compatible controller, signals and initial state (assumption; whether the data meet it is a Setup
%     or Results check: 4 kHz controller on recorded r - y versus recorded feedback force).
%     Main consequence: x^c_tau = 0 per window, no controller initialisation; consistent with the
%     encoder (model state and model controller both start on the machine's condition at tau). Contrast:
%     Kessels' direct form needs the machine controller state (Remark 5.4). Error before tau is not
%     carried in; the free-run validation does carry it.
% 14. Step order from no plant feedthrough: no algebraic loop, no added delay.
% 15. Every window sample scored (as Eq. 22); theta_base, theta_aug, eta joint, controller fixed.
%     Joint estimation framing (Dirk, 2026-10-06): the goal is to estimate theta_base jointly, in the ten
%     combinations. Obstacle from Jan's work: the split is not unique (2026 Sec. 5.2) and approximately
%     initialised parameters stay near their start (EJC Sec. 4.3, 2026 Sec. 6.3). Handover: Section IV
%     extends Jan's method with orthogonality so that the estimate of theta_base is UNIQUE (this
%     strength, not "recovers the true parameters"). State what interpretable means: the estimated
%     parameters keep their physical meaning; cite Gyorok 2026 Sec. 3 ("With non-unique theta_b, the
%     baseline model could lose its physical meaning"). Operational form only; the general definition
%     is the Introduction's (todo there). "Negation" and Kessels' definition (Sec. 5.3.1.3) belong to
%     Section IV; his Sec. 6.2 recommendation to the Introduction gap. Do NOT mention Jan's parameter regularisation anywhere.
%     Do NOT say explicitly that Section III's formulation is the unconstrained comparison arm; let
%     Section IV extend it naturally.
% Setup pattern for hyperparameters (2026 Sec. 7): honest basis per value (physical insight and
%     trial-and-error, line-search, inherited from earlier results).
% 16. Closed-loop fit can mask plant error; plant-domain checks and controller transfer in Results.
% Notation: u = P u_act as in Sections II and IV; no "(eom)" shorthand; controller matrices distinct from A_c(Y).
%
% CORRECTIONS TO THE BLOCK (cycle 10, skill step 2a; claims kept, form corrected)
% - Structure: five subsections. The skill (Order, one job per subsection) makes the closed-loop
%   simulation a component with its own subsection before the identification problem, and Known before
%   new puts the encoder after the simulation whose initial state it supplies. III-A model (item 1 to 7
%   except the non-uniqueness), III-B closed-loop simulation (items 12 to 14, new label sec:aug_loop),
%   III-C encoder (items 8 to 11, label sec:aug_encoder), III-D identification (item 15, 16, training
%   coordinates moved from II-C, optimiser, selection, M_U; label sec:aug_closed_loop kept because
%   Sections II and IV point there for the coordinates and the optimiser), III-E PS2 (new, label sec:aug_ps2,
%   which Section IV already references).
% - Item 7: the non-uniqueness of the split is a consequence of the identification problem, so it follows
%   eq:aug_loss in III-D (skill Order; block position yields to Known before new).
% - Item 9: W^a is Xavier uniform, gain 1 (D-239, thesis block wa_encoder_init='xavier_uniform_gain1'),
%   not Kaiming; the Kaiming todo is resolved.
% - Item 5: the DT correction after the step is the S-DP form of hoekstra2026lfr Table 1 (adopted), so
%   it needs no own reason; the todo is resolved by the citation.
% - Item 4: correcting every state row is the S-DP definition (hoekstra2025lfr Sec. 2: the learning
%   outputs affect the entire model state); X and Y authority D-103. Todo resolved, risk stated.
% - Item 12: Forssell cited without a section number (README source map: Automatica numbering unchecked).
% - Item 13: x^c_tau = 0 is a definition, not Kessels' Remark 5.4 reconstruction (code wins over D-142,
%   README open conflicts); the comparison with the encoder is dropped (encoder comes after III-B).
% - Notation (item "Notation"): the model input is u_act (normalised u^n), as the code feeds stage forces
%   and applies P inside f_base (FS/blocks.py::_deriv_with) and Section IV writes u^n = T_u(u_act - mu_u).
%   Encoder parameters theta_enc and training coordinates xi follow Section IV; the encoder history is
%   zeta (xi is taken), the machine controller state x^fb (s is the m_Delta scale of the appendix), the
%   encoder lag n_psi (superscript n means normalised), Toeplitz and input-to-state matrices calT, calC
%   (T_x and the reference r are taken), window starts calK.
%
% MUST ESTABLISH (PROPOSED, cycle 10, III-E, inferred from README ownership row III, source-map row
%   III-D, the README PS2 sketch of 2026-10-08 and DOC/PS2-METHOD.md)
%  1. Problem: the zero output layer gives the added states no gradient at the start; plain training
%     leaves most of the error of the omitted dynamics in place (THESIS-RESULTS step 2, pointer to VI).
%  2. S1 (display): runs alongside the minimisation of V, one update after every update of V; trains a
%     temporary head and the encoder's added-state part on the current model's closed-loop residual
%     (Hefny S1 adapted). S2 (display): one-step fit of the deployed g_aug on frozen encoder pairs
%     (Hefny S2 adapted). S3: plain minimisation of V continues (Downey Sec. 4.2). Switching rules as in
%     PS2Spec, values in Section V. No derivation.
%  3. M_PS2 named; THESIS-RESULTS step 3 terms.
%
% CYCLE 10 NOTES (header only)
% - Code confirmed this session: thesis block of ENT (_THESIS_FIXED, route all 6+n_a rows, Xavier W^a,
%   16 x 2 tanh, lag 29, closed loop, L-BFGS); GD/model.py::build_model (zero_init_feed_forward_nn,
%   Linear_Output_Block with Cd_norm, encoder from gantry_linearize_and_discretize at nominal parameters,
%   normalize_linear_ss_matrices); GD/model.py::get_encoder_dims (na_right = nb_right = 1, current
%   sample in the history); GD/data.py::compute_normalization (training records, q = P^-T y, velocity by
%   np.diff * fs, mean and std); FS/pre_encoder.py::linear_encoder_init_aug (W^b map, Xavier gain 1 on
%   one logical [W_y W_u], zero psi-tilde output layer, linear part on scaled uncentred histories,
%   x_off subtracted from the physical rows); FS/blocks.py::_rk4, _deriv_with (P inside, Horner, z, w),
%   Reduced_Gantry_State_Block.combinations_from_free (nine exp, m_diff bound with current m_total,
%   J_eff); FS/closed_loop.py (residual form, x^c = 0, step order, record-length-weighted validation
%   RMS); FS/interconnect.py loss (mse over batch, window and channels; param_loss 0 with the prior off)
%   and fit (after_step after every optimiser update); FS/ps2_opening.py (PS2Spec, after_step, _sample
%   pairs k, k+1 anywhere in the fit records, S2 normalised by the target variance, best calibration
%   weights restored); FS/lbfgs_polish.py::decide, meter_regression (tolerance 1e-3 relative);
%   GD/training.py::_gantry_meters (combo_err against the TRUE combinations, m_diff scaled by
%   0.5 (m1 + m2)). The param_loss meter is identically zero with the prior off and is not named.
% - Own claims checked by algebra this session: residual subtraction of eq:aug_direct_loop; centring
%   cancels in C^n (mu_y = P^T mu_q, positions of x_all are P^-T y exactly); encoder offset form
%   (W^b on scaled histories gives T_x x, minus T_x mu_x); zero-start gradient (x_bar reaches y_hat only
%   through f_aug, whose x_bar-derivative carries W_L of the physical rows); admissibility of
%   eq:mdelta_map for rho < 1; equilibria at Theta = 0 with u = 0 (K nonzero on Theta only);
%   observability of (A_d, C_d) from the three positions.
% - Resolved existing todos: missing bib entries (gyorok2026obc, forssell1999revisited added to refs.bib
%   with % CHECK lines; Forssell cited without section number); W^a Kaiming (D-239 Xavier); correction
%   after the RK4 step (adopted S-DP form); position rows (S-DP definition, D-103); open-loop drift
%   evidence (THESIS-RESULTS Sec. 5 plans it, pointer to Section VI); notation clash eta_0, s, xi_Delta
%   (encoder renamed theta_enc, zeta; controller state x^fb); nominal-parameter disclosure (stated).
%   Kept: VERIFY keys, VERIFY encoder bib, h_aug reason, LPV structure, zero-layer attribution,
%   Kessels Remark 5.3, figure labels, Y = 0, data-based initialisation, controller-rate check.
% - Left to Section V (its hyperparameter table, with the basis column): n_xbar, n_psi, n_f, stride, network
%   width and depth, Adam and L-BFGS settings, validation cadence, rho, the PS2 thresholds, and the window
%   choice from the closed-loop memory (D-220, header "Horizon check"); the masking check under a changed
%   controller to Section VI (THESIS-RESULTS step 6, E5).
% - Stale in other files (not edited): 04_obc.tex line 116 "Section sec:aug_structure showed that the
%   split ... is not unique" (now in sec:aug_closed_loop); line 159 "f^n_aug = T_x f_aug by eq:aug_norm"
%   (f_aug in physical units is no longer defined; f^n_aug is the first six outputs of phi_aug in
%   eq:aug_transition); line 189 "an open-loop simulation drifts (Section sec:aug_closed_loop)" (now
%   sec:aug_loop); eq:obc_model writes g_aug(theta_aug, x_hat_k, u_k) (now g_aug(theta_aug, x_hat^n_k,
%   u^n_k)). 05_setup.tex lines 55 and 56 (values and free-run validation "referenced from Section III-C",
%   now III-D) and the PS2 thresholds, which its table does not list. 99_appendix.tex app:lfr-details
%   still cites eq:lfr_chain (Section II's part). 02_system_baseline.tex uses "RK4" unexpanded in II-B,
%   before its expansion here.
% - Length: about 3.4 pages in a draftfalse build (scratchpad), against 1.55, so over 1.5 times after the
%   second deletion pass. Cause: the displays the Math standard requires (normalisation, model, zero start,
%   two loop displays, encoder and its initialisation, training coordinates, the problem display with all
%   constraint lines, two PS2 displays), the PS2 subsection and the coordinates moved from II-C, all
%   younger than the budget. Proposed budget: 3.3 pages including the two figures.
\section{Dynamic LPV-LFR Augmentation}\label{sec:augmentation}
\ifdraft
\begin{itemize}
\item Opening: from the baseline to a model that also captures omitted dynamics, identified under feedback; names III-A to III-E; delivers the second contribution of Section~I; one pointer to Section~V for records and values (style profile; Introduction)
\end{itemize}
\fi

Given the baseline of Section~II, our goal is a model that also
captures the dynamics it omits, identified from records measured under
feedback.
To this end, we define the augmented model
(Section~\ref{sec:aug_structure}), its simulation in closed loop with the
known controller (Section~\ref{sec:aug_loop}) and the encoder of its initial
state (Section~\ref{sec:aug_encoder}).
We then state the identification problem with joint estimation of
$\theta_{\mathrm{base}}$ (Section~\ref{sec:aug_closed_loop}) and a two-stage
predictive-state initialisation (PS2) of the added states
(Section~\ref{sec:aug_ps2}).
Together they deliver the second contribution of Section~I: dynamic-parallel
LFR augmentation extended to a closed-loop, self-scheduled state-space
setting with learned additional states.
Section~\ref{sec:setup} gives the records and the values of the symbols
introduced here.
\todo{Missing: contribution wording. This section also delivers joint estimation of $\theta_{\mathrm{base}}$ in the identifiable coordinates and the PS2 initialisation, which the second contribution does not name. Candidate: ``Extend dynamic-parallel LFR augmentation to a closed-loop, self-scheduled state-space setting, with joint estimation of the identifiable physical combinations and a data-driven initialisation of the learned additional states'' (D-222, D-229, D-236).}
\todo{VERIFY keys: \texttt{hoekstra2025lfr} = EJC 86 (2025) 101304,
\texttt{hoekstra2026lfr} = arXiv:2602.17297 (submitted to Automatica).
refs.bib and docs/references.md rows 46 and 130 now agree.}

\subsection{Augmented model}\label{sec:aug_structure}
\ifdraft
\begin{itemize}
\item Dynamic omitted effects need states that the six baseline states, a refit of $\theta_{\mathrm{base}}$ or a static augmentation cannot give (EJC Sec.~2; MUST ESTABLISH 1)
\item $f_{\mathrm{base}}$ is one RK4 step of the LFR of Section~II-B, input held; each stage schedules $\Dsch(Y)$ at its own state (D-018; this work)
\item Normalised coordinates from the training records, centred unlike Hoekstra Eq.~(28); centring cancels in the output map (code; TRAINING-DESIGN Sec.~3; (?) reason for centring)
\item Parallel, not series: the OBC of Section~IV is defined for additive augmentation; a series structure needs a residual network (Gy\"or\"ok 2026 abstract; EJC Sec.~4.3; D-003)
\item S-DP structure with a fixed interconnection, one tanh MLP for $f^{\mathrm n}_{\mathrm{aug}}$ and $g_{\mathrm{aug}}$ (EJC Sec.~3.2; 2026 Table~1, Sec.~6.3; D-180)
\item Input-output data fix $\bar x$ only up to a change of coordinates, so $\bar x$ has no physical meaning (MUST ESTABLISH 7; header Caution)
\item Zero output layer reproduces the baseline from the same physical initial state (2026 Eq.~(29); D-237)
\item The network corrects every state row, $X$ and $Y$ included, at the cost of accumulating position on the free axes; output map not augmented (?); no LPV structure (?); the only algebraic loop is the baseline's (EJC Sec.~2, Thm.~1; D-103; D-017)
\end{itemize}
\fi

Effects that the baseline omits need states of their own if they are
dynamic.
The six states of the baseline cannot represent them.
Refitting $\theta_{\mathrm{base}}$ adds no states.
A static augmentation corrects the state update of the baseline without
extending its state~\cite[Sec.~2]{hoekstra2025lfr}.
The augmentation therefore adds states $\bar x\in\R^{n_{\bar x}}$.

We realise the baseline block
$f_{\mathrm{base}}:\R^{10}\times\R^6\times\R^3\to\R^6$ as one fourth-order
Runge-Kutta (RK4) step of the LFR of Section~\ref{sec:lfr}
over the sample period $T_s$, with the actuator forces held over the step.
Each stage evaluates $\Dsch(Y)$ at its own state, so $f_{\mathrm{base}}$
remains self-scheduled.

The model acts on normalised signals, as in Hoekstra et
al.~\cite[Eq.~(28)]{hoekstra2026lfr}.
Unlike their scaling with state statistics from a simulation of the
baseline~\cite[Sec.~5.3]{hoekstra2026lfr}, we centre and scale with the
statistics of the training records:
\begin{equation}
\begin{aligned}
 \tilde x^{\mathrm n}&=T_x(\tilde x-\mu_x),\quad
 u^{\mathrm n}=T_u(\uact-\mu_u),\\
 y^{\mathrm n}&=T_y(y-\mu_y),\\
 f^{\mathrm n}_{\mathrm{base}}&(\theta_{\mathrm{base}},\tilde x^{\mathrm n},u^{\mathrm n})\\
 &=T_x\bigl(f_{\mathrm{base}}(\theta_{\mathrm{base}},
   T_x^{-1}\tilde x^{\mathrm n}+\mu_x,\\
 &\qquad\qquad T_u^{-1}u^{\mathrm n}+\mu_u)-\mu_x\bigr),\\
 C^{\mathrm n}&=T_y\begin{bmatrix}P^\top&0\end{bmatrix}T_x^{-1},
\end{aligned}
\label{eq:aug_norm}
\end{equation}
where $y\in\R^3$ is the recorded position $\qact$,
$\mu_x\in\R^6$, $\mu_u,\mu_y\in\R^3$ are means and
$T_x\in\R^{6\times6}$, $T_u,T_y\in\R^{3\times3}$ diagonal matrices of inverse
standard deviations over the training records, with the state
$\tilde x=\col(q,\dot q)$ formed from $q=P^{-\top}y$ and its difference
quotients, $f^{\mathrm n}_{\mathrm{base}}:\R^{10}\times\R^6\times\R^3\to\R^6$,
and $C^{\mathrm n}\in\R^{3\times6}$.
The centring cancels in the output map $C^{\mathrm n}$, because $\mu_y$ equals
$P^\top$ times the position part of $\mu_x$.
\todo{Missing: reason for centring and for record statistics instead of a baseline simulation. TRAINING-DESIGN Sec.~3 records only the user decision (2026-09-27). Candidate: an open-loop simulation of the baseline drifts on the free axes $X$ and $Y$, so its statistics would not describe the records (own reading, no decision entry).}

We augment the baseline in parallel, which keeps it an explicit term of the
state update.
We choose the parallel over a series interconnection because the
orthogonal-by-construction method that Section~\ref{sec:obc} extends is
defined for additive augmentation~\cite{gyorok2026obc}.
A series interconnection would also need a residual network to reproduce the
baseline at initialisation~\cite[Sec.~4.3]{hoekstra2025lfr}.

In the LFR augmentation framework of Hoekstra et al., a fixed interconnection
connects the baseline block with a learning
block~\cite[Sec.~3.2]{hoekstra2025lfr}.
We use its dynamic-parallel (S-DP) form~\cite[Table~1]{hoekstra2026lfr}, with
one tanh network as learning block as in its
experiments~\cite[Sec.~6.3]{hoekstra2026lfr} (Fig.~\ref{fig:aug_structure}),
\begin{equation}
\begin{aligned}
 \tilde x^{\mathrm n}_{k+1}
 &=f^{\mathrm n}_{\mathrm{base}}(\theta_{\mathrm{base}},\tilde x^{\mathrm n}_k,u^{\mathrm n}_k)
  +f^{\mathrm n}_{\mathrm{aug}}(\theta_{\mathrm{aug}},\hat x^{\mathrm n}_k,u^{\mathrm n}_k),\\
 \bar x_{k+1}&=g_{\mathrm{aug}}(\theta_{\mathrm{aug}},\hat x^{\mathrm n}_k,u^{\mathrm n}_k),\\
 \hat y^{\mathrm n}_k&=C^{\mathrm n}\tilde x^{\mathrm n}_k,\\
 \begin{bmatrix}
  f^{\mathrm n}_{\mathrm{aug}}\\ g_{\mathrm{aug}}
 \end{bmatrix}&(\theta_{\mathrm{aug}},\hat x^{\mathrm n},u^{\mathrm n})
 =\phi_{\mathrm{aug}}\bigl(\theta_{\mathrm{aug}},\col(\hat x^{\mathrm n},u^{\mathrm n})\bigr),
\end{aligned}
\label{eq:aug_transition}
\end{equation}
where $\hat x^{\mathrm n}=\col(\tilde x^{\mathrm n},\bar x)$,
$\hat y^{\mathrm n}\in\R^3$ is the normalised model output,
$\phi_{\mathrm{aug}}:\R^{n_{\mathrm{aug}}}\times\R^{9+n_{\bar x}}\to\R^{6+n_{\bar x}}$
is a multilayer perceptron with tanh hidden layers and a linear output layer,
$\theta_{\mathrm{aug}}\in\R^{n_{\mathrm{aug}}}$ are its parameters, and
$f^{\mathrm n}_{\mathrm{aug}}$ and $g_{\mathrm{aug}}$ are its first six and
last $n_{\bar x}$ outputs.
Input-output data fix $\bar x$ only up to an invertible change of
coordinates, so $\bar x$ has no physical meaning.
\todo{Missing: Section~\ref{sec:setup} introduces a truth-only absorber with $n_{\bar x}$ states; it must not identify $\bar x$ with the absorber states (header Caution; D-228: judge the added states at the output only). Candidate: one sentence there.}

\begin{figure}[t]
 \centering
 \resizebox{\columnwidth}{!}{\input{figures/tex/augmentation_structure}}
 \caption{Augmented model \eqref{eq:aug_transition}. The physics step
 $f_{\mathrm{base}}$ and the network $\phi_{\mathrm{aug}}$ act in parallel on
 $\hat x_k$; $\phi_{\mathrm{aug}}$ supplies $f^{\mathrm n}_{\mathrm{aug}}$ and
 $g_{\mathrm{aug}}$, and the output map is not learned
 ($h_{\mathrm{aug}}=0$). The encoder $\psi$ sets $\hat x$ at each window
 start. The input $u_k$ is the actuator force, during training the
 closed-loop input of Section~\ref{sec:aug_loop}.
 \todo{No setpoint generator currently for the inputs. to cramped together. We could make this a two column figure.}
 \todo{Figure labels: $u_k$ and $\hat x$ without the normalisation superscript of \eqref{eq:aug_transition}; the encoder histories end at $k-1$, while \eqref{eq:aug_encoder} includes sample $k$. The figure is generated and not edited here.}
 }
 \label{fig:aug_structure}
\end{figure}

Training starts from the baseline.
The output layer of $\phi_{\mathrm{aug}}$ is initialised at zero,
\begin{equation}
 W_L=0,\quad b_L=0
 \quad\Longrightarrow\quad
 f^{\mathrm n}_{\mathrm{aug}}=0,\quad g_{\mathrm{aug}}=0,
 \label{eq:aug_zero_init}
\end{equation}
where $W_L\in\R^{(6+n_{\bar x})\times n_h}$ and $b_L\in\R^{6+n_{\bar x}}$ are
its weight and bias and $n_h$ is the width of the last hidden layer.
From the same physical initial state, the augmented model then reproduces the
baseline~\cite[Eq.~(29)]{hoekstra2026lfr}.
\todo{Missing: attribution of the all-zero output layer. It follows the authors' public dynamic-parallel code (D-237); \cite[Sec.~5.4.3]{hoekstra2026lfr} also allows a random linear part. Ask Jan which produced the reported S-DP results.}

The network acts on every state row and through no other path.
Correcting the entire state is the S-DP form~\cite[Sec.~2]{hoekstra2025lfr}.
It gives the learned component authority on $X$ and $Y$, where omitted
effects can appear through the coupling of the axes in $M(Y)$.
On these free axes $K$ vanishes in \eqref{eq:damping_stiffness_matrices}, so
a persistent correction accumulates in position.
The output map is the kinematic relation of \eqref{eq:Ptransform}, which
Kessels leaves unaugmented when it is accurate~\cite[Remark~5.3]{kessels2025ai}.
The network depends on $Y$ through $\hat x^{\mathrm n}$, but unlike the
augmented LPV-LFR of Drenth~\cite[Eq.~(5.2)]{drenth2025thesis} it adds no
scheduling block.
No learned function feeds the baseline block within a step, so the
interconnection graph is acyclic~\cite[Thm.~1]{hoekstra2025lfr}.
The only algebraic loop is that of the baseline, whose well-posedness
\eqref{eq:wellposed} reduces to a nonsingular $M(Y)$.
\todo{Missing: reason for state augmentation only ($h_{\mathrm{aug}}=0$). Kessels' Remark~5.3 needs an accurate output map, which only Section~\ref{sec:setup} can establish for the data; D-140 item~5 rules out only an output augmentation that depends on the input, since it closes an algebraic loop with $D_{\mathrm{fb}}$~\cite[Remark~5.1]{kessels2025ai}. Candidate: the measured positions are the generalised coordinates through the exact map $P^\top$.}
\todo{Missing: reason for an unrestricted network rather than an LPV-structured augmentation. D-017 requires the augmentation in LFR $\Dsch(Y)$ form and no later decision supersedes it (README open conflicts); the research plan (Aspect~2) also proposed the LPV-LFR form. Candidate: Dirk decides whether D-017 is superseded.}

\subsection{Closed-loop simulation}\label{sec:aug_loop}
\ifdraft
\begin{itemize}
\item The criterion needs a simulation comparable with records measured under feedback; driving the model with the recorded input, as Hoekstra Eq.~(22), simulates it in open loop (meeting audit; D-140)
\item On the free axes an open-loop simulation leaves position offsets; the SUBNET consistency condition (Beintema Cond.~1, Hoekstra Sec.~5.5) fails in open loop and can hold in closed loop; Forssell's indirect approach as linear motivation; a motivation, not a proof; evidence in Section~VI (THESIS-RESULTS Sec.~5)
\item Machine and model loops share $r$, $u^{\mathrm{ff}}$ and the controller, in normalised coordinates (Kessels Eqs.~(5.13c), (5.13d))
\item Subtracting them gives the residual form with $x^{\mathrm c}_\tau=0$ (this work; D-140; code)
\item $x^{\mathrm c}_\tau=0$ places the model's controller in the machine's state; unlike Kessels Remark~5.4 no reconstruction (code)
\item No feedthrough of the model: fixed step order, no algebraic loop, no delay (D-140 item~5)
\item Controller at the model rate, cost reason (D-141); its measured deviation unreported (?)
\end{itemize}
\fi

The identification needs a simulation of the model that can be compared with
records measured under feedback.
The records hold the actuator forces and positions of the machine in closed
loop with a known controller.
Driving the model with the recorded forces, as the criterion of Hoekstra et
al.\ does~\cite[Eq.~(22)]{hoekstra2026lfr}, simulates it in open loop.
On the free axes, an error in the initial velocity of such a simulation
leaves a permanent position offset.
The consistency argument that Hoekstra et al.\ inherit from
SUBNET~\cite[Sec.~5.5]{hoekstra2026lfr} assumes an incrementally
exponentially output stable data-generating
system~\cite[Cond.~1]{beintema2023subnet}.
The open-loop gantry violates this condition, whereas its loop with the
stabilising controller can meet it.
For linear systems, fitting the response from the reference to the output
with a known controller is the indirect approach, in which an output-error
criterion remains consistent~\cite{forssell1999revisited}.
For our nonlinear model, these are motivations, not a proof of consistency.
Section~\ref{sec:results} compares the drift of open-loop and closed-loop
simulations on one record.

The machine and the model loop are driven by the reference $r$ and the
feedforward $u^{\mathrm{ff}}$ through the same linear time-invariant
controller (Fig.~\ref{fig:aug_closed_loop}).
They differ only in the output they feed back.

In normalised coordinates, the machine with controller state $x^{\mathrm{fb}}$
and the model with controller state $\hat x^{\mathrm{fb}}$ are
\begin{equation}
\begin{aligned}
 x^{\mathrm{fb}}_{k+1}&=A_{\mathrm{fb}}x^{\mathrm{fb}}_k+B^{\mathrm n}_{\mathrm{fb}}(r^{\mathrm n}_k-y^{\mathrm n}_k),\\
 u^{\mathrm n}_k&=C^{\mathrm n}_{\mathrm{fb}}x^{\mathrm{fb}}_k+D^{\mathrm n}_{\mathrm{fb}}(r^{\mathrm n}_k-y^{\mathrm n}_k)+u^{\mathrm{ff,n}}_k,\\
 \hat x^{\mathrm{fb}}_{k+1}&=A_{\mathrm{fb}}\hat x^{\mathrm{fb}}_k+B^{\mathrm n}_{\mathrm{fb}}(r^{\mathrm n}_k-\hat y^{\mathrm n}_k),\\
 \hat u^{\mathrm n}_k&=C^{\mathrm n}_{\mathrm{fb}}\hat x^{\mathrm{fb}}_k+D^{\mathrm n}_{\mathrm{fb}}(r^{\mathrm n}_k-\hat y^{\mathrm n}_k)+u^{\mathrm{ff,n}}_k,
\end{aligned}
\label{eq:aug_direct_loop}
\end{equation}
where $(A_{\mathrm{fb}},B_{\mathrm{fb}},C_{\mathrm{fb}},D_{\mathrm{fb}})$ is
the state-space controller from the position error $r-y$ to the actuator
forces, $x^{\mathrm{fb}},\hat x^{\mathrm{fb}}\in\R^{n_{\mathrm{fb}}}$,
$B^{\mathrm n}_{\mathrm{fb}}=B_{\mathrm{fb}}T_y^{-1}$,
$C^{\mathrm n}_{\mathrm{fb}}=T_uC_{\mathrm{fb}}$ and
$D^{\mathrm n}_{\mathrm{fb}}=T_uD_{\mathrm{fb}}T_y^{-1}$ are its input and
output matrices in normalised coordinates, $r^{\mathrm n}=T_y(r-\mu_y)$,
$u^{\mathrm{ff,n}}=T_u(u^{\mathrm{ff}}-\mu_u)$, and $\hat u^{\mathrm n}$ is the
model input that replaces $u^{\mathrm n}$ in \eqref{eq:aug_transition}.
The model loop has the feedback equations of Kessels' closed-loop
criterion~\cite[Eqs.~(5.13c), (5.13d)]{kessels2025ai}, with the error formed
against the model output.

\begin{figure}[t]
 \centering
 \includegraphics[width=\columnwidth]{closed_loop_same_reference.pdf}
 \caption{Recorded plant loop (top) and augmented-model loop (bottom),
 driven by one shared reference $r_k$ and feedforward $u_k^{\mathrm{ff}}$,
 with the same feedback controller. The model output is evaluated
 before the input advances its state to the next sample.}
 \label{fig:aug_closed_loop}
\end{figure}

We simulate the model loop in a residual form, obtained by subtracting the
machine's equations from the model's:
\begin{equation}
\begin{aligned}
 x^{\mathrm c}_{k+1}&=A_{\mathrm{fb}}x^{\mathrm c}_k+B^{\mathrm n}_{\mathrm{fb}}(y^{\mathrm n}_k-\hat y^{\mathrm n}_k),
 \qquad x^{\mathrm c}_\tau=0,\\
 \hat u^{\mathrm n}_k&=u^{\mathrm n}_k+C^{\mathrm n}_{\mathrm{fb}}x^{\mathrm c}_k
 +D^{\mathrm n}_{\mathrm{fb}}(y^{\mathrm n}_k-\hat y^{\mathrm n}_k),
\end{aligned}
\label{eq:aug_residual_loop}
\end{equation}
where $x^{\mathrm c}=\hat x^{\mathrm{fb}}-x^{\mathrm{fb}}$ and $\tau$ is the
first sample of a simulated window.
The model thus receives the recorded input plus the controller's response to
its own output error, without $r$ or $u^{\mathrm{ff}}$.

The residual form equals the model loop of \eqref{eq:aug_direct_loop} when the
recorded input was produced by this controller from the recorded positions.
\todo{Check for Setup: the training loop uses the controller re-discretised at the model rate, while the records were generated at a higher rate. Run the re-discretised controller on the recorded $r-y$ and compare its feedback force with the recorded one. A rollout with $\hat y=y$ gives $\hat u=u$ by construction and checks nothing.}

Because $x^{\mathrm c}$ is a difference of controller states,
$x^{\mathrm c}_\tau=0$ places the model's controller in the machine's
controller state at every window start.
Model error made before $\tau$ therefore does not enter the window.
Kessels' direct form instead needs the machine's controller state at every
window start, read from the machine or reconstructed from the reference, the
output and the controller~\cite[Remark~5.4]{kessels2025ai}.

The model has no feedthrough from $\hat u^{\mathrm n}$ to
$\hat y^{\mathrm n}$, which fixes the order of each sample.
First $\hat y^{\mathrm n}_k$ follows from $\tilde x^{\mathrm n}_k$, then
$\hat u^{\mathrm n}_k$ from \eqref{eq:aug_residual_loop}, and then the model
and controller states advance.
The controller feedthrough $D^{\mathrm n}_{\mathrm{fb}}$ therefore closes no
algebraic loop.
The loop also adds no delay.

The controller runs at the model sampling rate, whereas the records were
generated at a higher rate.
Running the loop at the record rate would multiply the cost of every
simulation by the ratio of the rates.
\todo{Missing: the deviation that re-discretisation causes is measured (D-141: the peak output sensitivity near 150~Hz moves by 15.3\,\%, the phase margin by 3.6~degrees) but no section reports it, and D-141 requires it to be stated wherever closed-loop numbers are compared with the records. Candidate: one sentence in Section~\ref{sec:setup} next to the controller rates.}

\subsection{Encoder and initial state}\label{sec:aug_encoder}
\ifdraft
\begin{itemize}
\item Each window needs an initial state, of which only positions are recorded; an encoder estimates it from the history (Beintema Sec.~2 to 3)
\item Linear map plus nonlinear correction (encoder paper Eq.~(8)); the linear part acts on scaled, uncentred histories, the correction on normalised ones (code, D-055)
\item The encoder estimates all states, the measured positions included (?: Kessels Remark~5.3)
\item The encoder paper assumes a stable baseline; ours is marginally stable but observable from the positions (encoder paper Sec.~2; this work)
\item $W^b$ from the reconstructability map of the ZOH linearisation at rest at $Y=0$ with the nominal parameters; $W^a$ Xavier; zero output layer of $\tilde\psi$ (encoder paper Eqs.~(10) to (17), (30) to (32); 2026 Eq.~(31); D-239)
\item Computed, not fitted as 2026 Eq.~(30); better start and faster convergence at comparable accuracy (encoder paper Sec.~4.4); nominal parameters while $\theta_{\mathrm{base}}$ starts elsewhere; co-estimation outside the paper's scope (footnote~1)
\item Open: why $Y=0$; model-based versus data-based (?)
\end{itemize}
\fi

Each simulated window needs an initial state $\hat x^{\mathrm n}_\tau$, of
which the records hold only the positions.
Following Beintema et al., an encoder estimates it from the recorded
history~\cite{beintema2023subnet}.
As in the encoder of Hoekstra et
al.~\cite[Eq.~(8)]{hoekstra2026encoder}, a linear map is combined with a
nonlinear correction,
\begin{equation}
 \hat x^{\mathrm n}_{\tau|\tau}
 =\psi(\theta_{\mathrm{enc}},\zeta^{\mathrm n}_\tau)
 =\begin{bmatrix}W^b\\ W^a\end{bmatrix}(\zeta^{\mathrm n}_\tau+\mu_\zeta)
 -\begin{bmatrix}T_x\mu_x\\ 0\end{bmatrix}
 +\tilde\psi(\zeta^{\mathrm n}_\tau),
 \label{eq:aug_encoder}
\end{equation}
where $\zeta^{\mathrm n}_\tau\in\R^{6(n_\psi+1)}$ stacks $y^{\mathrm n}$ and
then $u^{\mathrm n}$ over the samples $\tau-n_\psi,\dots,\tau$,
$\mu_\zeta=\col(\mathbf 1_{n_\psi+1}\otimes T_y\mu_y,\,
\mathbf 1_{n_\psi+1}\otimes T_u\mu_u)$ restores their means,
$W^b\in\R^{6\times6(n_\psi+1)}$ and $W^a\in\R^{n_{\bar x}\times6(n_\psi+1)}$
give the physical and the added states,
$\tilde\psi:\R^{6(n_\psi+1)}\to\R^{6+n_{\bar x}}$ is a multilayer perceptron,
and $\theta_{\mathrm{enc}}\in\R^{n_{\mathrm{enc}}}$ collects $W^b$, $W^a$ and
the parameters of $\tilde\psi$.
The offsets $\mu_\zeta$ and $T_x\mu_x$ let $W^b$ act on uncentred signals, as
the map of a linearisation about the origin requires.
\todo{Missing: whether to state that the encoder also estimates the measured positions, and to mention the alternative of initialising them from the measurement~\cite[Remark~5.3]{kessels2025ai}.}
\todo{VERIFY bib: \texttt{hoekstra2026encoder} = arXiv:2602.13108v1
(13 Feb 2026), Hoekstra, Gy\"or\"ok, T\'oth, Schoukens; matches PDF p.~1.
Check for a published (IFAC) version before submission.}

The physical rows start from the baseline, following the model-based
initialisation of Hoekstra et al.~\cite[Sec.~3.1]{hoekstra2026encoder}.
That initialisation assumes a stable baseline~\cite[Sec.~2]{hoekstra2026encoder},
while ours is only marginally stable.

For zero input, every rest state with $\Theta=0$ is an equilibrium of
\eqref{eq:eom}.
Linearised at the one with $Y=0$ and discretised by zero-order hold, the
baseline gives the initial encoder
\begin{equation}
\begin{aligned}
 W^b&=\begin{bmatrix}
  A_d^{n_\psi}\mathcal O^{+} & \;\mathcal C-A_d^{n_\psi}\mathcal O^{+}\mathcal T
 \end{bmatrix},\\
 W^a_{ij}&\sim\mathcal U\Bigl(-\sqrt{6/(n_{\bar x}+6n_\psi+6)},\,
   \sqrt{6/(n_{\bar x}+6n_\psi+6)}\Bigr),\\
 \tilde\psi&\text{: output layer }0,
\end{aligned}
\label{eq:aug_enc_init}
\end{equation}
where $(A_d,B_d,C_d,D_d)$ is that discretisation over $T_s$, with the nominal
parameters of Appendix~\ref{app:params} and in the scaled coordinates of
\eqref{eq:aug_norm}, $\mathcal O\in\R^{3(n_\psi+1)\times6}$ is its
observability matrix over $n_\psi+1$ samples, $\mathcal T\in\R^{3(n_\psi+1)\times3(n_\psi+1)}$
the Toeplitz matrix of its Markov parameters, and
$\mathcal C\in\R^{6\times3(n_\psi+1)}$ its input-to-state map over the same
samples~\cite[Eqs.~(10) to (17)]{hoekstra2026encoder}.
The measured positions make $(A_d,C_d)$ observable, so $W^b$ exists despite
the free axes.

The encoder starts from the baseline's linearisation and a random
added-state map and is trained with the model.
The Xavier draw of $W^a$ follows~\cite[Eq.~(31)]{hoekstra2026lfr}.
Unlike~\cite[Eq.~(30)]{hoekstra2026lfr}, which fits the baseline encoder to
simulated baseline states, we compute $W^b$ directly, since the linearised
baseline is known.
This model-based initialisation gives better starting values and faster
convergence than a random encoder, at comparable final
accuracy~\cite[Sec.~4.4]{hoekstra2026encoder}.
$W^b$ uses the nominal parameters also when $\theta_{\mathrm{base}}$ starts
elsewhere (Section~\ref{sec:aug_closed_loop}).
The encoder paper leaves the co-estimation of $\theta_{\mathrm{base}}$ out of
scope~\cite[footnote~1]{hoekstra2026encoder}.
\todo{Missing: why linearise at $Y=0$. The encoder paper also linearises ``around the 0 point''~\cite[Sec.~4.2]{hoekstra2026encoder} but gives no reason, and the $Y$-scheduled $W^b(Y)$ of D-167 is not wired.}
\todo{Missing: why the model-based rather than the data-based initialisation~\cite[Sec.~3.2]{hoekstra2026encoder}. Candidate: the data-based fit needs forward simulations of the baseline on the data, which drift on the free axes (own reading); or negligible computation~\cite[Sec.~4.4]{hoekstra2026encoder}.}

\subsection{Closed-loop identification}\label{sec:aug_closed_loop}
\ifdraft
\begin{itemize}
\item $\theta_{\mathrm{base}}$ is estimated jointly in coordinates that keep \eqref{eq:lfr_admissibility}: nine logarithmic, $m_\Delta$ bounded (D-229; code); they bound the inertia, not the component masses
\item $\vartheta=(\xi,\theta_{\mathrm{aug}},\theta_{\mathrm{enc}})$ trained jointly, unlike Kessels Eq.~(5.12); controller and $L_b$ fixed; start $\xi=0$ (D-222)
\item Criterion: every sample of truncated closed-loop windows from encoder states, as 2026 Eq.~(22) (Beintema Sec.~2 to 3; D-220)
\item Adam, then an L-BFGS polish as in Drenth (D-171)
\item Selection: closed-loop free run on the validation records; the polish kept only if it lowers it and the combination error, against the true combinations, does not rise (D-171, D-172; code; oracle (?))
\item The split is not unique (2026 Sec.~5.2, 6.3); interpretable means $\theta_{\mathrm{base}}$ keeps its physical meaning, which needs a unique estimate; Section~IV (Gy\"or\"ok 2026 Sec.~3; MUST ESTABLISH 15)
\item A closed-loop fit can mask plant error; Section~VI evaluates under a changed controller (header Caution; THESIS-RESULTS step~6)
\item $\mathcal M_{\mathrm U}$ named, inputs and output, purpose in THESIS-RESULTS terms (README ownership; notation (?))
\end{itemize}
\fi

The model must now be fitted to the records, with $\theta_{\mathrm{base}}$
estimated jointly in coordinates that keep every iterate admissible
(Section~\ref{sec:params}).
We train the combinations through the coordinates $\xi\in\R^{10}$, which
bound $m_\Delta$ but not the individual masses $m_1$ and $m_2$:
\begin{equation}
\begin{aligned}
 \theta_{\mathrm{base},j}(\xi)&=\theta_{\mathrm{base},0,j}\,e^{\xi_j},
 \qquad \theta_{\mathrm{base},j}\neq m_\Delta,\\
 m_\Delta(\xi)&=\rho\,\frac{2}{L_b}\sqrt{m_\Sigma(\xi)J_{\mathrm{eff}}(\xi)}\,
   \tanh(\eta_0+s\,\xi_\Delta),
\end{aligned}
\label{eq:mdelta_map}
\end{equation}
where $\theta_{\mathrm{base},0}\in\R^{10}$ is the start, $\xi_\Delta$ the
entry of $\xi$ that belongs to $m_\Delta$, $\rho\in(0,1)$ a margin, and
$\eta_0,s\in\R$ make $\xi=0$ reproduce $\theta_{\mathrm{base},0}$ with
$\partial m_\Delta/\partial\xi_\Delta=m_{\Delta,0}$
(Appendix~\ref{app:lfr-wellposed}).
Since $\rho<1$, every $\xi$ gives
$\tfrac{L_b^2}{4}m_\Delta^2<m_\Sigma J_{\mathrm{eff}}$, which is
\eqref{eq:lfr_admissibility}.

Unlike Kessels, who keeps the first-principles parameters
fixed~\cite[Eq.~(5.12)]{kessels2025ai}, we train
$\vartheta=(\xi,\theta_{\mathrm{aug}},\theta_{\mathrm{enc}})
\in\R^{10+n_{\mathrm{aug}}+n_{\mathrm{enc}}}$ jointly.
The controller and $L_b$ stay fixed.
Training starts at $\xi=0$, the network of \eqref{eq:aug_zero_init} and the
encoder of \eqref{eq:aug_enc_init}.

Truncated windows keep the cost of an update independent of the record length
and stabilise the optimisation~\cite[Secs.~2, 3]{beintema2023subnet}.
As in~\cite[Eq.~(22)]{hoekstra2026lfr}, the criterion therefore scores every
sample of such windows, simulated from encoder states:
\begin{equation}
\begin{aligned}
 &\min_{\vartheta}\;V(\vartheta)=\frac{1}{|\mathcal K|\,n_f\,n_y}
 \sum_{\tau\in\mathcal K}\sum_{\ell=0}^{n_f-1}
 \norm{y^{\mathrm n}_{\tau+\ell}-\hat y^{\mathrm n}_{\tau+\ell|\tau}}_2^2
 \quad\text{s.t.}\\
 &\tilde x^{\mathrm n}_{k+1}
 =f^{\mathrm n}_{\mathrm{base}}\bigl(\theta_{\mathrm{base}}(\xi),\tilde x^{\mathrm n}_k,\hat u^{\mathrm n}_k\bigr)
  +f^{\mathrm n}_{\mathrm{aug}}(\theta_{\mathrm{aug}},\hat x^{\mathrm n}_k,\hat u^{\mathrm n}_k),\\
 &\bar x_{k+1}=g_{\mathrm{aug}}(\theta_{\mathrm{aug}},\hat x^{\mathrm n}_k,\hat u^{\mathrm n}_k),
 \qquad \hat y^{\mathrm n}_{k|\tau}=C^{\mathrm n}\tilde x^{\mathrm n}_k,\\
 &\hat u^{\mathrm n}_k=u^{\mathrm n}_k+C^{\mathrm n}_{\mathrm{fb}}x^{\mathrm c}_k
 +D^{\mathrm n}_{\mathrm{fb}}(y^{\mathrm n}_k-\hat y^{\mathrm n}_{k|\tau}),\\
 &x^{\mathrm c}_{k+1}=A_{\mathrm{fb}}x^{\mathrm c}_k
 +B^{\mathrm n}_{\mathrm{fb}}(y^{\mathrm n}_k-\hat y^{\mathrm n}_{k|\tau}),
 \qquad x^{\mathrm c}_\tau=0,\\
 &\hat x^{\mathrm n}_\tau=\psi(\theta_{\mathrm{enc}},\zeta^{\mathrm n}_\tau),
\end{aligned}
\label{eq:aug_loss}
\end{equation}
where $\mathcal K$ is a set of window starts in the training records,
$n_f$ the window length, $n_y=3$, and the constraints hold for
$k=\tau,\dots,\tau+n_f-1$.
Through $T_y$, $V$ weights each position channel with its inverse variance.

Adam minimises $V$ over minibatches of windows for a fixed number of updates.
As in Drenth et al.~\cite[Sec.~5]{drenth2025lpvlfr}, an L-BFGS polish on a
fixed set of windows follows.

The parameters are selected on a closed-loop free run of the validation
records.
At a fixed interval of updates, the model runs in the loop
\eqref{eq:aug_residual_loop} over each whole validation record, from the
encoder state at its first complete history and $x^{\mathrm c}=0$.
The selected Adam iterate minimises the mean over the records of the RMS
position error in metres, weighted by record length.
The polish is kept only if it lowers this validation error.
It must also not raise the combination error by more than a relative
tolerance.
The combination error is the RMS of the errors of the ten combinations
relative to their true values.
The error of $m_\Delta$ is divided by the mean mass of the two X drives
instead.
\todo{Missing: the polish acceptance uses the true combinations (\texttt{combo\_err} in \texttt{GD/training.py::\_gantry\_meters}), an oracle inside the training procedure (sources.md trap 6). Candidate: drop the meter from the acceptance rule for the reported runs, or report it as oracle-gated; Dirk decides.}

From the zero start, $f^{\mathrm n}_{\mathrm{aug}}$ can also learn changes that
a different $\theta_{\mathrm{base}}$ would produce, so the minimiser of
\eqref{eq:aug_loss} is not unique in the split between baseline and
network~\cite[Sec.~5.2]{hoekstra2026lfr}.
In the experiments of Hoekstra et al., approximately initialised physical
parameters stay close to their start, while the learning components learn
dynamics that the baseline could represent~\cite[Sec.~6.3]{hoekstra2026lfr}.
With a non-unique $\theta_{\mathrm{base}}$, the baseline can lose its physical
meaning~\cite[Sec.~3]{gyorok2026obc}.
We call the estimated combinations interpretable if they keep this meaning,
which requires a unique estimate.
Section~\ref{sec:obc} extends the method so that this estimate is unique.

A good closed-loop fit does not show that the model is correct in open loop,
because the controller acts on the output error and can mask plant error.
Section~\ref{sec:results} therefore also evaluates the final model under a
changed controller.

We call the model \eqref{eq:aug_transition} identified by \eqref{eq:aug_loss}
from \eqref{eq:aug_zero_init} and $\theta_{\mathrm{base},0}$, and selected as
above, $\mathcal M_{\mathrm U}$.
Its inputs are the recorded actuator forces and positions.
Its output is the predicted position $\hat y$.
Section~\ref{sec:results} uses $\mathcal M_{\mathrm U}$ to show what plain
training from the zero start achieves.
\todo{Missing: model notation. Candidate: $\mathcal M_{\mathrm U}$ and $\mathcal M_{\mathrm{PS2}}$, calligraphic as $\mathcal M_{\mathrm{LPV}}$, since $M$ is the inertia matrix (README open conflict on model names).}

\subsection{Added-state initialisation}\label{sec:aug_ps2}
\ifdraft
\begin{itemize}
\item With the zero output layer, $\bar x$ does not reach the output, so $V$ gives $W^a$ and the added-state rows of the output layer no gradient at the start; plain training leaves most of the error of the omitted dynamics in place (PS2-METHOD Sec.~2.2; THESIS-RESULTS step~2; this work)
\item PS2 (this work, adapting Hefny Sec.~3): S1 alongside $V$, one update after every update of $V$, trains a temporary network and the added-state part of the encoder on the current closed-loop residual (D-234, D-236; code)
\item S1 ends on a plateau or cap of the unexplained fraction on withheld training records; the screen skips S2 (PS2Spec; README sketch)
\item S2 fits $g_{\mathrm{aug}}$ once to propagate the encoder's added states one step (Hefny Sec.~3 adapted; code: pairs $k$, $k+1$ anywhere in the fit records)
\item S2 stops on a plateau or cap of its normalised criterion on withheld records; the best parameters are kept (PS2Spec)
\item S3 continues the minimisation of $V$ unchanged (Downey Sec.~4.2)
\item $\mathcal M_{\mathrm{PS2}}$ named, inputs and output, purpose in THESIS-RESULTS step~3 terms (notation (?))
\end{itemize}
\fi

With the zero output layer of \eqref{eq:aug_zero_init}, $\bar x$ does not
reach the output, so $V$ gives $W^a$ and the added-state rows of $W_L$ no
gradient at the start.
From this start, plain training lowers the error below the baseline but
leaves most of the error of the omitted dynamics in place
(Section~\ref{sec:results}).
We therefore initialise the added-state part from data with PS2, in three
stages that run inside the minimisation of $V$.

S1 runs alongside that minimisation, with one S1 update after every update of
$V$.
Like the first-stage regression of Hefny et
al.~\cite[Sec.~3]{hefny2015supervised}, it regresses future signals on the
history, here the current model's closed-loop residuals through the added
states:
\begin{equation}
\begin{aligned}
 &\min_{\omega,\,\theta^a_{\mathrm{enc}}}\;
 \sum_{\tau\in\mathcal K_1}
 \norm{\kappa(\omega,\bar x_{\tau|\tau})-e_\tau/\sigma_e}_2^2,\\
 &e_\tau=\col\bigl(y^{\mathrm n}_{\tau+\ell}
 -\hat y^{\mathrm n}_{\tau+\ell|\tau}\bigr)_{\ell=0}^{n_\psi},
\end{aligned}
\label{eq:ps2_s1}
\end{equation}
where $e_\tau\in\R^{3(n_\psi+1)}$, $\bar x_{\tau|\tau}$ is the added-state part
of \eqref{eq:aug_encoder},
$\hat y^{\mathrm n}_{\tau+\ell|\tau}$ the output of the current model in the
loop \eqref{eq:aug_residual_loop},
$\kappa:\R^{n_\omega}\times\R^{n_{\bar x}}\to\R^{3(n_\psi+1)}$ a temporary
network with one tanh hidden layer and parameters $\omega\in\R^{n_\omega}$,
$\theta^a_{\mathrm{enc}}$ the encoder parameters that produce $\bar x$ ($W^a$
and the added-state rows of the output layer of $\tilde\psi$), $\sigma_e>0$
the RMS of the targets $e_\tau$, and $\mathcal K_1$ a set of window starts
drawn for every update from the training records that are not withheld below.
S1 thus trains the encoder to put into $\bar x$ the information that predicts
the current residual.

At a fixed interval of updates, $\kappa$ is evaluated on fixed windows of
training records withheld from S1 and S2.
The evaluated quantity is the unexplained fraction $\gamma$, the ratio of the
sum of $\norm{\kappa(\omega,\bar x_{\tau|\tau})-e_\tau/\sigma_e}_2^2$ to the
sum of $\norm{e_\tau/\sigma_e}_2^2$ over these windows.
S1 ends after a minimum number of updates once the lowest $\gamma$ of the last
four evaluations is no longer a set fraction below the lowest $\gamma$ before
them.
It ends at the latest after a maximum number of updates.
If $\gamma$ exceeds a set level at a fixed update or at the end of S1, S2 is
skipped.

At the end of S1, S2 fits the deployed added-state transition once, like the
second-stage regression of~\cite[Sec.~3]{hefny2015supervised} but over the
nonlinear $g_{\mathrm{aug}}$ instead of a linear operator:
\begin{equation}
 \min_{\theta_{\mathrm{aug}}}\;
 \sum_{k\in\mathcal P}
 \norm{g_{\mathrm{aug}}\bigl(\theta_{\mathrm{aug}},\hat x^{\mathrm n}_{k|k},u^{\mathrm n}_k\bigr)
 -\bar x_{k+1|k+1}}_2^2,
 \label{eq:ps2_s2}
\end{equation}
where $\hat x^{\mathrm n}_{k|k}$ and $\bar x_{k+1|k+1}$ are outputs of
\eqref{eq:aug_encoder} at $k$ and $k+1$ with $\theta_{\mathrm{enc}}$ fixed at
the end of S1, and $\mathcal P$ is a set of samples drawn from the training
records that are not withheld.

At a fixed interval of steps, the criterion of \eqref{eq:ps2_s2} is evaluated
on pairs from the withheld records, divided by the variance of their targets
$\bar x_{k+1|k+1}$.
S2 stops after a minimum number of steps once the lowest such value of the
last five evaluations is no longer a set fraction below the lowest value
before them.
It stops at the latest after a maximum number of steps.
S2 keeps the network parameters with the lowest evaluated value.

S3 discards $\kappa$ and continues the minimisation of $V$ from the modified
$\theta_{\mathrm{enc}}$ and $\theta_{\mathrm{aug}}$, as Downey et al.\
refine a two-stage regression initialisation by backpropagation through
time~\cite[Sec.~4.2]{downey2017psrnn}.
The model \eqref{eq:aug_transition}, the criterion \eqref{eq:aug_loss}, the
optimiser and the selection stay unchanged.
\todo{Missing: Section~\ref{sec:setup} lists no values for the PS2 rules (evaluation interval, minimum and maximum updates and steps, plateau fractions, screen level and update, number of withheld windows and pairs, batch of $\mathcal K_1$; \texttt{PS2Spec}, all HEURISTIC). Candidate: rows of its hyperparameter table with basis ``heuristic'' (PS2-METHOD Sec.~5).}

We call the model \eqref{eq:aug_transition} identified by \eqref{eq:aug_loss}
from $\theta_{\mathrm{base},0}$ with PS2, and selected as in
Section~\ref{sec:aug_closed_loop}, $\mathcal M_{\mathrm{PS2}}$.
Its inputs and output are those of $\mathcal M_{\mathrm U}$.
Section~\ref{sec:results} compares $\mathcal M_{\mathrm{PS2}}$ with
$\mathcal M_{\mathrm U}$ to show whether, with PS2, the added states learn the
omitted dynamics.
